\nochapter{致谢}

在本次课程设计的资料查阅、需求梳理、数据库脚本联调与文档撰写过程中，得到了多方面的支持和帮助。

首先，感谢 \textbf{王文玉} 老师的耐心指导。老师在课程设计的各个阶段都给予了宝贵的建议，特别是在 E-R 模型设计、规范化处理和存储过程编写等方面提供了专业的指导，使本设计能够更加规范和完整。

其次，感谢同学们在环境配置、测试数据验证和系统联调方面的帮助。在数据库初始化、API 测试和前端功能验证过程中，同学们提出了许多宝贵的意见和建议，帮助我发现并修复了多个问题。

再次，感谢学校图书馆和网络资源的支持。通过查阅《数据库系统概论》、MySQL 官方文档等资料，使我能够深入理解数据库设计的理论基础，并将其应用到实践中。

最后，感谢所有为这个项目做出贡献的人。虽然这是一个课程设计项目，但它体现了团队合作的精神，也让我学到了很多实用的技能和知识。

在此，我对所有给予帮助的人表示衷心的感谢。因水平与时间所限，本设计中难免存在疏漏之处，恳请老师和同学们批评指正。

\vspace{2cm}

\begin{flushright}
  学生签名：李昊阳  \\
  日期：2026.5.9 
\end{flushright}
